%% kapa: exam=1 round=27 year=2016 period=2 form=A subject=appraisal_law subject_ko=감정평가관계법규 seq=2 number=2
%% source: Q-Net 감정평가사 기출문제 (cst00309) artlSeq=5169402 file=extracted/1008/5169402/제27회 감정평가사 1차 문제/제27회 감정평가사 2교시문제(A형).pdf
%% conversion: kordoc 4.24.51 -> tools/gen1.py (auto); review=proofread-sonnet; flags=none
\begin{question}{2}
국토의 계획 및 이용에 관한 법령상 기반시설의 구분과 시설 종류의 연결이 옳지 않은 것은?
\begin{choices}
  \item 공간시설: 공원, 운동장
  \item 유통ㆍ공급시설: 유통업무설비, 방송ㆍ통신시설
  \item 보건위생시설: 화장시설, 도축장
  \item 환경기초시설: 하수도, 폐차장
  \item 방재시설: 하천, 유수지
\end{choices}
\end{question}
\answer{1}
